\documentclass[12pt,a4paper]{article}
\usepackage[utf8]{inputenc}
\usepackage[brazil]{babel}
\usepackage{graphicx}
\usepackage{amssymb, amsfonts, amsmath}
\usepackage{float}
\usepackage{enumerate}
\usepackage[top=2.5cm, bottom=2.5cm, left=1.25cm, right=1.25cm]{geometry}

\begin{document}
\pagestyle{empty}

\begin{center}
  \begin{tabular}{ccc}
    \begin{tabular}{c}
      \includegraphics[scale=0.25]{../../biblioteca/imagem/brasao-de-armas-brasil} \\
    \end{tabular} & 
    \begin{tabular}{c}
      Ministério da Educação \\
      Universidade Federal dos Vales do Jequitinhonha e Mucuri \\
      Faculdade de Ciências Sociais, Aplicadas e Exatas - FACSAE \\
      Departamento de Ciências Exatas - DCEX \\
      Disciplina: Introdução à Lógica \quad Semestre: 2022/1\\
      Prof. Dr. Luiz C. M. de Aquino\\
    \end{tabular} &
    \begin{tabular}{c}
      \includegraphics[scale=0.25]{../../biblioteca/imagem/logo-ufvjm} \\
    \end{tabular}
  \end{tabular}
\end{center}

\begin{center}
  \textbf{Lista I}
\end{center}

\begin{enumerate}
  \item (FGV/TCE-SE -- Adaptada) Considere a afirmação: ``Se hoje é sábado, então amanhã não trabalharei''. 
  A negação dessa afirmação é:

    \begin{enumerate}
      \item Hoje é sábado e amanhã trabalharei.
      \item Hoje não é sábado e amanhã trabalharei.
      \item Hoje não é sábado ou amanhã trabalharei.
      \item Se hoje não é sábado, então amanhã trabalharei.
      \item Se hoje não é sábado, então amanhã não trabalharei.
    \end{enumerate}

  \item (FGV – TJSC – 2015) Considere a sentença: ``Se cometi um crime, então
    serei condenado''. Uma sentença logicamente equivalente à sentença dada é:

    \begin{enumerate}
      \item  Não cometi um crime ou serei condenado.
      \item  Se não cometi um crime, então não serei condenado.
      \item  Se eu for condenado, então cometi um crime.
      \item  Cometi um crime e serei condenado.
      \item  Não cometi um crime e não serei condenado.
    \end{enumerate}

  \item (CESPE – TRE/MT – 2015) A negação da proposição: ``Se o número
    inteiro $m > 2$ é primo, então o número $m$ é ímpar'' pode ser expressa
    corretamente por:

    \begin{enumerate}
      \item ``Se o número $m$ não é ímpar, então o número inteiro $m > 2$ não é primo''.
      \item ``Se o número inteiro $m > 2$ não é primo, então o número $m$ é ímpar''.
      \item ``O número inteiro $m > 2$ é primo e o número $m$ não é ímpar''.
      \item ``O número inteiro $m > 2$ é não primo e o número $m$ é ímpar''.
      \item ``Se o número inteiro $m > 2$ não é primo, então o número $m$ não é ímpar''
    \end{enumerate}

  \item Quatro suspeitos de praticar um crime fazem as seguintes declarações:
  
    \begin{itemize}
      \item João: Carlos é o criminoso;
      \item Pedro: eu não sou criminoso;
      \item Carlos: Paulo é o criminoso;
      \item Paulo: Carlos está mentindo;
    \end{itemize}

  Sabendo que APENAS UM dos suspeitos mente, determine quem é o criminoso.

    \begin{enumerate}
      \item João.
      \item Pedro.
      \item Carlos.
      \item Paulo.
    \end{enumerate}

  \item (FCC – SABESP – 2014) Alan, Beto, Caio e Décio são irmãos e foram
    interrogados pela própria mãe para saber quem comeu, sem autorização,
    o chocolate que estava no armário. Sabe-se que apenas um dos quatro
    comeu o chocolate, e que os quatro irmãos sabem quem foi. A mãe
    perguntou para cada um quem cometeu o ato, ao que recebeu as seguintes
    respostas:

    \begin{itemize}   
      \item Alan diz que foi Beto;
      \item Beto diz que foi Caio;
      \item Caio diz que Beto mente;
      \item Décio diz que não foi ele.
    \end{itemize}

    O irmão que fala a verdade e o irmão que comeu o chocolate são,
    respectivamente,

    \begin{enumerate}
      \item Beto e Décio.
      \item Alan e Beto.
      \item Beto e Caio.
      \item Alan e Caio.
      \item Caio e Décio
    \end{enumerate}
        
        
    \item (Vunesp/TJ-SP - Adaptada) Em um edifício com apartamentos somente nos andares de 1º ao 4º, moram 4 meninas, 
    em andares distintos: Juliana, Carla, Alice e Maria, não necessariamente nessa ordem. 
    Cada uma delas tem um animal de estimação diferente: gato, cachorro, passarinho e tartaruga, 
    não necessariamente nessa ordem. Maria vive reclamando do barulho feito pelo cachorro, 
    no andar imediatamente acima do seu. Juliana, que não mora no 4º andar, mora um andar acima do de Alice, 
    que tem o passarinho e não mora no 2º andar. Quem mora no 3º andar tem uma tartaruga. 
    Determine em qual andar mora cada menina e qual seu respectivo animal de estimação.

\end{enumerate}

\begin{center}
  \textbf{Gabarito}
\end{center}

[1] (a). 
[2] (a).
[3] (c).
[4] (c). 
[5] 
Alice mora no 1º andar e tem um passarinho. 
Juliana mora no 2º andar e tem um gato. 
Maria mora no 3º andar e tem um tartaruga.  
Carla mora no 4º andar e tem um cachorro.

\end{document}